

\section{Цель работы}

Целью данной лабораторной работы является освоение постановки и решения
транспортной задачи с промежуточными складами. В отличие от классической
транспортной задачи, в которой груз перевозится непосредственно от пунктов
отправления к пунктам назначения, здесь маршрут включает два последовательных
этапа: сначала груз доставляется из пунктов отправления $A_i$ на склады $D_j$,
а затем со складов --- в пункты назначения $B_k$. Каждый склад имеет
ограниченную ёмкость, поэтому задача не сводится к двум независимым
транспортным задачам: потоки первого и второго этапов связаны условием
сохранения груза на складе.

В рамках работы требуется:
\begin{itemize}
    \item освоить содержательную постановку двухэтапной транспортной задачи
          и условия её совместности;
    \item построить математическую модель задачи в виде задачи линейного
          программирования с ограничениями на запасы, заявки и ёмкости
          складов;
    \item научиться представлять задачу в виде сети и сводить её к задаче
          о потоке минимальной стоимости;
    \item освоить метод последовательных кратчайших путей с потенциалами
          для нахождения потока минимальной стоимости;
    \item приобрести практические навыки проверки оптимальности найденного
          плана с помощью остаточных стоимостей;
    \item научиться интерпретировать результат: какие объёмы груза и по каким
          маршрутам нужно перевезти, чтобы суммарная стоимость была
          минимальной.
\end{itemize}

В конечном счёте выполнение работы должно сформировать практические навыки
применения методов дискретной оптимизации и сетевых алгоритмов к задачам
планирования перевозок.

\clearpage

\section{Теоретическая часть}

\subsection*{Постановка задачи}

Имеются $m$ пунктов отправления $A_1,\ldots,A_m$, $s$ промежуточных складов
$D_1,\ldots,D_s$ и $n$ пунктов назначения $B_1,\ldots,B_n$. Запасы груза в
пунктах отправления равны $a_i$, потребности пунктов назначения --- $b_k$,
ёмкости складов --- $d_j$. Стоимость перевозки единицы груза из $A_i$ в $D_j$
равна $c^{AD}_{ij}$, а из $D_j$ в $B_k$ --- $c^{DB}_{jk}$. Требуется найти
объёмы перевозок $x_{ij}$ (из $A_i$ в $D_j$) и $y_{jk}$ (из $D_j$ в $B_k$),
при которых суммарная стоимость всех перевозок минимальна.

Содержательно задача описывает ситуацию, когда прямой транзит груза от
поставщиков к потребителям невозможен или невыгоден, а распределение груза
происходит через сеть складов. Склады не производят и не потребляют груз:
сколько груза в склад привезено, столько из него и должно быть вывезено.
Именно это условие связывает два этапа перевозок в единую задачу.

\subsection*{Условия совместности}

Для существования допустимого плана нужно, чтобы выполнялись два условия.
Во-первых, суммарный запас должен быть равен суммарной потребности:
\[
\sum_{i=1}^{m} a_i = \sum_{k=1}^{n} b_k.
\]
Если это равенство не выполняется, задача балансируется введением
фиктивного поставщика или фиктивного потребителя с нулевыми стоимостями.
Во-вторых, суммарная ёмкость складов должна быть достаточной для
пропуска всего груза:
\[
\sum_{j=1}^{s} d_j \ge \sum_{i=1}^{m} a_i.
\]
В данной работе оба условия выполнены: суммарный запас и суммарная
потребность равны 92, а суммарная ёмкость складов равна 103.

\subsection*{Математическая модель}

Введём неотрицательные переменные $x_{ij}$ и $y_{jk}$, означающие
количество груза, перевозимого по соответствующим маршрутам. Задача
принимает вид
\[
\min Z=\sum_{i=1}^{m}\sum_{j=1}^{s} c^{AD}_{ij}x_{ij}
      +\sum_{j=1}^{s}\sum_{k=1}^{n} c^{DB}_{jk}y_{jk}
\]
при ограничениях
\begin{align}
&\sum_{j=1}^{s} x_{ij}=a_i, && i=1,\ldots,m, \label{eq:supply}\\
&\sum_{i=1}^{m} x_{ij}=\sum_{k=1}^{n} y_{jk}, && j=1,\ldots,s, \label{eq:balance}\\
&\sum_{j=1}^{s} y_{jk}=b_k, && k=1,\ldots,n, \label{eq:demand}\\
&\sum_{i=1}^{m} x_{ij}\le d_j, && j=1,\ldots,s, \label{eq:capacity}\\
&x_{ij}\ge 0,\quad y_{jk}\ge 0. \label{eq:nonneg}
\end{align}

Ограничения \eqref{eq:supply} требуют вывезти весь запас каждого пункта
отправления, ограничения \eqref{eq:demand} --- удовлетворить все заявки,
ограничения \eqref{eq:capacity} не позволяют превысить ёмкость склада, а
равенства \eqref{eq:balance} выражают сохранение груза на складе. Целевая
функция суммирует стоимости перевозок на обоих этапах.

Матрица ограничений этой задачи является матрицей сети, поэтому она
обладает свойством полной унимодулярности. Это означает, что при целых
исходных данных среди оптимальных решений задачи линейного
программирования существует целочисленное. Поэтому нет необходимости
отдельно накладывать требование целочисленности переменных: достаточно
решить непрерывную задачу, и оптимальный план можно выбрать целым.

\subsection*{Сетевая интерпретация}

Задачу удобно представить в виде ориентированной сети с одним источником
$S$ и одним стоком $T$. Вершинами сети являются:
\begin{itemize}
    \item источник $S$ и сток $T$;
    \item пункты отправления $A_1,\ldots,A_m$;
    \item входы складов $D_j^{\mathrm{in}}$ и их выходы
          $D_j^{\mathrm{out}}$, $j=1,\ldots,s$;
    \item пункты назначения $B_1,\ldots,B_n$.
\end{itemize}

Дуги сети и их параметры приведены в таблице~\ref{tab:network-arcs}.

\begin{table}[h!]
\centering
\caption{Дуги сети и их параметры}
\label{tab:network-arcs}
\begin{tabular}{|c|c|c|c|}
\hline
Дуга & Пропускная способность & Стоимость & Смысл \\
\hline
$S\to A_i$ & $a_i$ & 0 & запас пункта $A_i$ \\
\hline
$A_i\to D_j^{\mathrm{in}}$ & $+\infty$ & $c^{AD}_{ij}$ & перевозка $A_i\to D_j$ \\
\hline
$D_j^{\mathrm{in}}\to D_j^{\mathrm{out}}$ & $d_j$ & 0 & ёмкость склада $D_j$ \\
\hline
$D_j^{\mathrm{out}}\to B_k$ & $+\infty$ & $c^{DB}_{jk}$ & перевозка $D_j\to B_k$ \\
\hline
$B_k\to T$ & $b_k$ & 0 & потребность пункта $B_k$ \\
\hline
\end{tabular}
\end{table}

Схематически сеть имеет вид
\[
S \xrightarrow{\ a_i\ } A_i
  \xrightarrow{\ c^{AD}_{ij}\ } D_j^{\mathrm{in}}
  \xrightarrow{\ d_j\ } D_j^{\mathrm{out}}
  \xrightarrow{\ c^{DB}_{jk}\ } B_k
  \xrightarrow{\ b_k\ } T .
\]

Поток величины $F=\sum_i a_i=\sum_k b_k$ в этой сети соответствует
допустимому плану перевозок. Действительно, поток по дуге
$S\to A_i$ не может превысить $a_i$, а поскольку суммарный поток равен
$\sum_i a_i$, каждая такая дуга должна быть насыщена. Аналогично насыщаются
дуги $B_k\to T$, что обеспечивает выполнение всех заявок. Дуга
$D_j^{\mathrm{in}}\to D_j^{\mathrm{out}}$ ограничивает суммарный груз,
проходящий через склад, а равенство входа и выхода склада следует из
сохранения потока в вершинах $D_j^{\mathrm{in}}$ и $D_j^{\mathrm{out}}$.
Если поток содержит циркуляционные циклы, их можно удалить, не изменив
величину потока и лишь уменьшив стоимость, поэтому оптимальный поток
циклов не содержит и взаимно однозначно определяет план перевозок.

Таким образом, исходная задача эквивалентна задаче о потоке минимальной
стоимости: требуется пропустить $F$ единиц потока из $S$ в $T$ с
минимальными суммарными затратами.

\subsection*{Остаточная сеть и условие оптимальности}

Пусть $f_e$ --- поток по дуге $e=(u,v)$ с пропускной способностью $u_e$ и
стоимостью $c_e$. Остаточной пропускной способностью называется величина
$u_e-f_e$ в прямом направлении и $f_e$ в обратном. Остаточная сеть содержит
для каждой дуги две остаточные дуги: прямую с пропускной способностью
$u_e-f_e$ и стоимостью $c_e$ и обратную с пропускной способностью $f_e$ и
стоимостью $-c_e$. Обратная дуга означает возможность уменьшить ранее
назначенный поток, то есть перераспределить груз между маршрутами.

Поток $f$ величины $F$ является потоком минимальной стоимости тогда и
только тогда, когда в остаточной сети нет отрицательного цикла. Если бы
такой цикл существовал, вдоль него можно было бы увеличить поток на
небольшую величину и уменьшить суммарную стоимость, не изменив величину
потока $F$. Эквивалентное условие: среди всех потоков данной величины нет
потока меньшей стоимости, то есть не существует отрицательного по стоимости
пути из $S$ в $T$ в остаточной сети.

Для проверки этого условия используются потенциалы вершин. Функция
$h:V\to\mathbb{R}$ называется потенциалом, если для каждой остаточной дуги
$(u,v)$ выполняется неравенство
\[
c^{h}_{uv}=c_{uv}+h(u)-h(v)\ge 0.
\]
Величина $c^{h}_{uv}$ называется остаточной (приведённой) стоимостью.
Если потенциалы найдены, то вместо исходных стоимостей можно использовать
неотрицательные приведённые стоимости, и для поиска кратчайшего пути
становится применим алгоритм Дейкстры.

\subsection*{Метод последовательных кратчайших путей}

Метод последовательных кратчайших путей (successive shortest paths)
строит оптимальный поток постепенно, увеличивая его вдоль кратчайших путей
в остаточной сети. Алгоритм работает следующим образом.

\begin{enumerate}
    \item Положить $f_e=0$ для всех дуг и $h(v)=0$ для всех вершин.
    \item Пока величина потока меньше $F$, выполнять следующие шаги:
    \begin{enumerate}
        \item найти кратчайший путь из $S$ в $T$ в остаточной сети по
              приведённым стоимостям $c^{h}_{uv}$ с помощью алгоритма
              Дейкстры;
        \item если путь не существует, то поток величины $F$ невозможен;
        \item для всех достижимых вершин обновить потенциалы:
              $h(v)\leftarrow h(v)+\operatorname{dist}(v)$, где
              $\operatorname{dist}(v)$ --- кратчайшее расстояние от $S$
              до $v$;
        \item вычислить пропускную способность пути
              $\Delta=\min_{e\in P}(u_e-f_e)$ и увеличить поток вдоль
              пути на $\Delta$;
        \item увеличить суммарную стоимость на
              $\Delta\cdot\sum_{e\in P}c_e$.
    \end{enumerate}
    \item Выдать поток $f$ и его стоимость.
\end{enumerate}

Обновление потенциалов гарантирует, что после каждого шага все приведённые
стоимости остаточных дуг неотрицательны. Поэтому на следующей итерации
снова можно применять алгоритм Дейкстры, несмотря на наличие отрицательных
стоимостей у обратных дуг. Если все пропускные способности целые, то
величина $\Delta$ также целая, поэтому алгоритм строит целочисленный поток.

Трудоёмкость метода составляет $O(F\,m\log n)$ при использовании
двоичной кучи, где $F$ --- величина потока, $m$ --- число дуг, $n$ ---
число вершин. В данной задаче $F=92$, $m=83$, $n=25$, поэтому решение
находится за доли секунды.

\subsection*{Связь с методом потенциалов}

Классический метод потенциалов для транспортной задачи также опирается на
двойственные переменные $u_i$ и $v_j$, для которых выполняются неравенства
$u_i+v_j\le c_{ij}$. Разность $c_{ij}-(u_i+v_j)$ играет роль приведённой
стоимости: она показывает, насколько перевозка по данной клетке дороже
оптимальной оценки. В сетевой постановке потенциалы вершин $h(v)$ являются
обобщением потенциалов $u_i$ и $v_j$, а условие
$c_{uv}+h(u)-h(v)\ge 0$ --- обобщением условия оптимальности метода
потенциалов. Таким образом, метод последовательных кратчайших путей можно
рассматривать как сетевой вариант метода потенциалов, приспособленный для
задач с промежуточными складами и ограничениями на пропускные способности.

\clearpage

\section{Практическая часть}

\subsection*{Исходные данные}

В соответствии с вариантом 4 заданы шесть пунктов отправления
$A_1,\ldots,A_6$, пять складов $D_1,\ldots,D_5$ и семь пунктов назначения
$B_1,\ldots,B_7$. Запасы груза приведены в таблице~\ref{tab:supply},
потребности --- в таблице~\ref{tab:demand}, ёмкости складов --- в
таблице~\ref{tab:capacity}.

\begin{table}[h!]
\centering
\caption{Запасы груза в пунктах отправления}
\label{tab:supply}
\begin{tabular}{|c|c|c|c|c|c|c|}
\hline
Пункт & $A_1$ & $A_2$ & $A_3$ & $A_4$ & $A_5$ & $A_6$ \\
\hline
Запас $a_i$ & 17 & 13 & 13 & 12 & 18 & 19 \\
\hline
\end{tabular}
\end{table}

\begin{table}[h!]
\centering
\caption{Заявки пунктов назначения}
\label{tab:demand}
\begin{tabular}{|c|c|c|c|c|c|c|c|}
\hline
Пункт & $B_1$ & $B_2$ & $B_3$ & $B_4$ & $B_5$ & $B_6$ & $B_7$ \\
\hline
Заявка $b_k$ & 10 & 14 & 13 & 10 & 12 & 11 & 22 \\
\hline
\end{tabular}
\end{table}

\begin{table}[h!]
\centering
\caption{Ёмкости складов}
\label{tab:capacity}
\begin{tabular}{|c|c|c|c|c|c|}
\hline
Склад & $D_1$ & $D_2$ & $D_3$ & $D_4$ & $D_5$ \\
\hline
Ёмкость $d_j$ & 23 & 17 & 20 & 21 & 22 \\
\hline
\end{tabular}
\end{table}

Суммарный запас равен
\[
17+13+13+12+18+19=92,
\]
суммарная потребность также равна
\[
10+14+13+10+12+11+22=92,
\]
а суммарная ёмкость складов составляет
\[
23+17+20+21+22=103.
\]
Следовательно, условия совместности выполнены: суммарный запас равен
суммарной потребности, а ёмкости складов достаточно для пропуска всего
груза.

Стоимости перевозок из пунктов отправления на склады приведены в
таблице~\ref{tab:cad}, а со складов в пункты назначения --- в
таблице~\ref{tab:cdb}.

\begin{table}[h!]
\centering
\caption{Стоимости перевозок $c^{AD}_{ij}$ из $A_i$ в $D_j$}
\label{tab:cad}
\begin{tabular}{|c|c|c|c|c|c|}
\hline
 & $D_1$ & $D_2$ & $D_3$ & $D_4$ & $D_5$ \\
\hline
$A_1$ & 6 & 6 & 6 & 7 & 9 \\
\hline
$A_2$ & 7 & 5 & 7 & 2 & 9 \\
\hline
$A_3$ & 6 & 3 & 7 & 2 & 6 \\
\hline
$A_4$ & 9 & 8 & 4 & 8 & 9 \\
\hline
$A_5$ & 6 & 2 & 5 & 8 & 3 \\
\hline
$A_6$ & 4 & 9 & 6 & 8 & 8 \\
\hline
\end{tabular}
\end{table}

\begin{table}[h!]
\centering
\caption{Стоимости перевозок $c^{DB}_{jk}$ из $D_j$ в $B_k$}
\label{tab:cdb}
\small
\begin{tabular}{|c|c|c|c|c|c|c|c|}
\hline
 & $B_1$ & $B_2$ & $B_3$ & $B_4$ & $B_5$ & $B_6$ & $B_7$ \\
\hline
$D_1$ & 6 & 5 & 4 & 6 & 7 & 2 & 6 \\
\hline
$D_2$ & 4 & 4 & 5 & 3 & 4 & 3 & 4 \\
\hline
$D_3$ & 4 & 3 & 2 & 4 & 4 & 6 & 7 \\
\hline
$D_4$ & 2 & 4 & 6 & 6 & 4 & 5 & 5 \\
\hline
$D_5$ & 7 & 3 & 8 & 6 & 3 & 4 & 3 \\
\hline
\end{tabular}
\end{table}

\subsection*{Построение сети}

Построим сеть с источником $S$ и стоком $T$. Она содержит
$1+6+5+5+7+1=25$ вершин. Число дуг равно
\[
6+6\cdot 5+5+5\cdot 7+7=6+30+5+35+7=83.
\]
Пропускные способности дуг $S\to A_i$ равны запасам $a_i$, дуг
$B_k\to T$ --- заявкам $b_k$, а дуг $D_j^{\mathrm{in}}\to D_j^{\mathrm{out}}$
--- ёмкостям $d_j$. Стоимости дуг $A_i\to D_j^{\mathrm{in}}$ и
$D_j^{\mathrm{out}}\to B_k$ равны соответственно $c^{AD}_{ij}$ и
$c^{DB}_{jk}$. Требуемая величина потока равна $F=92$.

\subsection*{Ход решения}

Решение выполнялось методом последовательных кратчайших путей с
потенциалами. На каждой итерации в остаточной сети находился кратчайший
путь из $S$ в $T$, после чего поток вдоль этого пути увеличивался на
максимально возможную величину. В таблице~\ref{tab:trace} приведена трасса
вычислений: номер итерации, дополнительный поток $\Delta$, стоимость
единицы потока на выбранном пути, накопленная величина потока и
накопленная стоимость.

\begin{table}[h!]
\centering
\caption{Трасса метода последовательных кратчайших путей}
\label{tab:trace}
\footnotesize
\begin{tabular}{|c|c|c|c|c|p{6.7cm}|}
\hline
№ & $\Delta$ & $c(P)$ & $f$ & $Z$ & Маршрут (с учётом остаточных дуг) \\
\hline
1  & 10 & 4  & 10 & 40  & $A_2\to D_4\to B_1$ \\
\hline
2  & 10 & 5  & 20 & 90  & $A_5\to D_2\to B_4$ \\
\hline
3  & 7  & 5  & 27 & 125 & $A_5\to D_2\to B_6$ \\
\hline
4  & 3  & 6  & 30 & 143 & $A_2\to D_4\to B_2$ \\
\hline
5  & 8  & 6  & 38 & 191 & $A_3\to D_4\to B_2$ \\
\hline
6  & 1  & 6  & 39 & 197 & $A_5\to D_5\to B_2$ \\
\hline
7  & 12 & 6  & 51 & 269 & $A_4\to D_3\to B_3$ \\
\hline
8  & 4  & 6  & 55 & 293 & $A_6\to D_1\to B_6$ \\
\hline
9  & 2  & 7  & 57 & 307 & $A_6\to D_1\to B_6$, затем $D_2\to B_2$ \\
\hline
10 & 5  & 7  & 62 & 342 & $A_3\to D_2$, вытеснение $A_5\to D_2$, далее $A_5\to D_5\to B_5$ \\
\hline
11 & 5  & 7  & 67 & 377 & $A_6\to D_1\to B_6$, затем $D_2\to B_5$ \\
\hline
12 & 1  & 8  & 68 & 385 & $A_6\to D_1\to B_3$ \\
\hline
13 & 2  & 9  & 70 & 403 & $A_6\to D_1\to B_2$, затем $D_4\to B_5$ \\
\hline
14 & 1  & 9  & 71 & 412 & $A_6\to D_1\to B_2$, затем $D_5\to B_7$ \\
\hline
15 & 4  & 9  & 75 & 448 & $A_6\to D_1\to B_2$, цепочка $D_4\to B_5$, $D_2\to B_7$ \\
\hline
16 & 3  & 9  & 78 & 475 & $A_1\to D_3\to B_2$, цепочка $D_4\to B_5$, $D_2\to B_7$ \\
\hline
17 & 4  & 9  & 82 & 511 & $A_1\to D_3\to B_2$, цепочка $D_4\to B_5$, $D_5\to B_7$ \\
\hline
18 & 1  & 10 & 83 & 521 & $A_1\to D_3\to B_2$, замена $D_1\to B_2$ на $D_1\to B_7$ \\
\hline
19 & 9  & 10 & 92 & 611 & $A_1\to D_2$, вытеснение $A_5\to D_2$, далее $A_5\to D_5\to B_7$ \\
\hline
\end{tabular}
\end{table}

На первых восьми итерациях все выбранные пути являются простыми путями
$A_i\to D_j\to B_k$ без обратных дуг. Начиная с девятой итерации алгоритм
использует обратные дуги остаточной сети: это означает перераспределение
ранее назначенных перевозок. В таблице~\ref{tab:trace} слово <<затем>>
означает переход по обратной дуге, то есть уменьшение ранее назначенного
потока на соответствующем участке. Например, на итерации 10 груз из $A_3$
направляется на склад $D_2$, а ранее назначенный поток $A_5\to D_2$
уменьшается на 5 единиц; освободившийся груз пункта $A_5$ направляется по
маршруту $A_5\to D_5\to B_5$. Такие перераспределения не увеличивают
величину потока, но уменьшают его стоимость, поэтому метод последовательных
кратчайших путей продолжает работу до тех пор, пока не будет достигнута
величина $F=92$ и в остаточной сети не останется отрицательных путей.

После 19 итераций получен поток величины 92 и минимальной стоимостью 611.
Матрица перевозок $X=(x_{ij})$ из пунктов отправления на склады приведена
в таблице~\ref{tab:X}, а матрица перевозок $Y=(y_{jk})$ со складов в пункты
назначения --- в таблице~\ref{tab:Y}.

\begin{table}[h!]
\centering
\caption{Матрица перевозок $X$ из $A_i$ в $D_j$}
\label{tab:X}
\begin{tabular}{|c|c|c|c|c|c|c|}
\hline
 & $D_1$ & $D_2$ & $D_3$ & $D_4$ & $D_5$ & Итого \\
\hline
$A_1$ & 0 & 9 & 8 & 0 & 0 & 17 \\
\hline
$A_2$ & 0 & 0 & 0 & 13 & 0 & 13 \\
\hline
$A_3$ & 0 & 5 & 0 & 8 & 0 & 13 \\
\hline
$A_4$ & 0 & 0 & 12 & 0 & 0 & 12 \\
\hline
$A_5$ & 0 & 3 & 0 & 0 & 15 & 18 \\
\hline
$A_6$ & 19 & 0 & 0 & 0 & 0 & 19 \\
\hline
Итого & 19 & 17 & 20 & 21 & 15 & 92 \\
\hline
\end{tabular}
\end{table}

\begin{table}[h!]
\centering
\caption{Матрица перевозок $Y$ из $D_j$ в $B_k$}
\label{tab:Y}
\small
\begin{tabular}{|c|c|c|c|c|c|c|c|c|}
\hline
 & $B_1$ & $B_2$ & $B_3$ & $B_4$ & $B_5$ & $B_6$ & $B_7$ & Итого \\
\hline
$D_1$ & 0 & 6 & 1 & 0 & 0 & 11 & 1 & 19 \\
\hline
$D_2$ & 0 & 0 & 0 & 10 & 0 & 0 & 7 & 17 \\
\hline
$D_3$ & 0 & 8 & 12 & 0 & 0 & 0 & 0 & 20 \\
\hline
$D_4$ & 10 & 0 & 0 & 0 & 11 & 0 & 0 & 21 \\
\hline
$D_5$ & 0 & 0 & 0 & 0 & 1 & 0 & 14 & 15 \\
\hline
Итого & 10 & 14 & 13 & 10 & 12 & 11 & 22 & 92 \\
\hline
\end{tabular}
\end{table}

\clearpage

\subsection*{Проверка допустимости плана}

Проверим выполнение ограничений задачи. Суммы элементов строк матрицы $X$
совпадают с запасами:
\[
17,\quad 13,\quad 13,\quad 12,\quad 18,\quad 19.
\]
Суммы элементов столбцов матрицы $Y$ совпадают с заявками:
\[
10,\quad 14,\quad 13,\quad 10,\quad 12,\quad 11,\quad 22.
\]
Суммы элементов строк матрицы $Y$ равны суммам элементов столбцов матрицы
$X$, то есть количество груза, привезённого на каждый склад, равно
количеству груза, вывезенного с него. Наконец, суммарный груз,
прошедший через каждый склад, не превышает его ёмкости
(таблица~\ref{tab:warehouse-load}).

\begin{table}[h!]
\centering
\caption{Загрузка складов и их ёмкости}
\label{tab:warehouse-load}
\begin{tabular}{|c|c|c|c|}
\hline
Склад & Привезено, ед. & Ёмкость, ед. & Свободный остаток \\
\hline
$D_1$ & 19 & 23 & 4 \\
\hline
$D_2$ & 17 & 17 & 0 \\
\hline
$D_3$ & 20 & 20 & 0 \\
\hline
$D_4$ & 21 & 21 & 0 \\
\hline
$D_5$ & 15 & 22 & 7 \\
\hline
\end{tabular}
\end{table}

Склады $D_2$, $D_3$ и $D_4$ загружены полностью, а склады $D_1$ и $D_5$
имеют свободный остаток. Все ограничения по запасам, заявкам, сохранению
груза и ёмкостям выполнены.

\subsection*{Стоимость перевозок}

Вычислим стоимость перевозок по каждому этапу. Для маршрутов $A_i\to D_j$
имеем:
\[
\begin{aligned}
Z_{AD}={}&9\cdot 6+8\cdot 6+13\cdot 2+5\cdot 3+8\cdot 2\\
&+12\cdot 4+3\cdot 2+15\cdot 3+19\cdot 4\\
={}&54+48+26+15+16+48+6+45+76=334.
\end{aligned}
\]
Для маршрутов $D_j\to B_k$:
\[
\begin{aligned}
Z_{DB}={}&6\cdot 5+1\cdot 4+11\cdot 2+1\cdot 6+10\cdot 3+7\cdot 4\\
&+8\cdot 3+12\cdot 2+10\cdot 2+11\cdot 4+1\cdot 3+14\cdot 3\\
={}&30+4+22+6+30+28+24+24+20+44+3+42=277.
\end{aligned}
\]
Суммарная стоимость:
\[
Z_{\min}=Z_{AD}+Z_{DB}=334+277=611.
\]

Распределение стоимости по складам приведено в
таблице~\ref{tab:cost-by-warehouse}. Видно, что наибольший вклад даёт
склад $D_3$ (144 единицы), а наименьший --- склад $D_5$ (90 единиц).

\begin{table}[h!]
\centering
\caption{Стоимость перевозок по складам}
\label{tab:cost-by-warehouse}
\begin{tabular}{|c|c|c|c|}
\hline
Склад & Стоимость $A\to D$ & Стоимость $D\to B$ & Итого \\
\hline
$D_1$ & 76 & 62 & 138 \\
\hline
$D_2$ & 75 & 58 & 133 \\
\hline
$D_3$ & 96 & 48 & 144 \\
\hline
$D_4$ & 42 & 64 & 106 \\
\hline
$D_5$ & 45 & 45 & 90 \\
\hline
Итого & 334 & 277 & 611 \\
\hline
\end{tabular}
\end{table}

\subsection*{Проверка оптимальности}

Для проверки оптимальности найдём потенциалы вершин остаточной сети.
В качестве потенциалов можно взять кратчайшие расстояния от источника $S$
в остаточной сети, вычисленные на последней итерации алгоритма. Значения
потенциалов приведены ниже:
\[
\begin{array}{c|cccccc}
 & A_1&A_2&A_3&A_4&A_5&A_6\\ \hline
h & 0&3&3&2&4&0
\end{array}
\]
\[
\begin{array}{c|ccccc}
 & D_1&D_2&D_3&D_4&D_5\\ \hline
h^{\mathrm{in}}  & 4&6&6&5&7\\
h^{\mathrm{out}} & 4&6&6&6&7
\end{array}
\]
\[
\begin{array}{c|ccccccc}
 & B_1&B_2&B_3&B_4&B_5&B_6&B_7\\ \hline
h & 8&9&8&9&10&6&10
\end{array}
\]
Кроме того, $h(S)=0$ и $h(T)=10$.

Для всех дуг остаточной сети приведённые стоимости
$c^{h}_{uv}=c_{uv}+h(u)-h(v)$ неотрицательны. В частности, для использованных
дуг $A_i\to D_j$ и $D_j\to B_k$ они равны нулю:
\[
\begin{aligned}
c^{h}(A_1,D_2)&=6+0-6=0, &
c^{h}(A_2,D_4)&=2+3-5=0,\\
c^{h}(A_5,D_2)&=2+4-6=0, &
c^{h}(D_1,B_6)&=2+4-6=0.
\end{aligned}
\]
Для незагруженных дуг приведённые стоимости также неотрицательны,
например:
\[
c^{h}(A_1,D_1)=6+0-4=2,\qquad
c^{h}(A_2,D_2)=5+3-6=2,
\]
\[
c^{h}(D_1,B_1)=6+4-8=2,\qquad
c^{h}(D_5,B_1)=7+7-8=6.
\]
Наконец, для насыщенных дуг $D_j^{\mathrm{in}}\to D_j^{\mathrm{out}}$
прямая остаточная дуга отсутствует, а обратная дуга имеет неотрицательную
приведённую стоимость. Например, для склада $D_4$
\[
c^{h}(D_4^{\mathrm{out}},D_4^{\mathrm{in}})
=0+6-5=1\ge 0.
\]
Таким образом, в остаточной сети нет дуг с отрицательной приведённой
стоимостью, а значит, нет и отрицательного цикла. Следовательно, найденный
поток является потоком минимальной стоимости, а план перевозок ---
оптимальным.

\subsection*{Интерпретация результата}

Оптимальный план имеет следующие особенности. Пункт $A_6$ направляет весь
свой запас (19 единиц) на склад $D_1$; пункт $A_2$ полностью использует
склад $D_4$ (13 единиц); пункт $A_4$ целиком направляет груз на склад $D_3$
(12 единиц). Пункты $A_1$, $A_3$ и $A_5$ распределяют груз между двумя
складами: $A_1$ --- между $D_2$ и $D_3$, $A_3$ --- между $D_2$ и $D_4$,
$A_5$ --- между $D_2$ и $D_5$.

Со складов груз распределяется следующим образом. Склад $D_1$ обслуживает
пункты $B_2$, $B_3$, $B_6$ и $B_7$, причём основная часть (11 единиц)
направляется в $B_6$. Склад $D_2$ обеспечивает $B_4$ и $B_7$, склад $D_3$
--- $B_2$ и $B_3$, склад $D_4$ --- $B_1$ и $B_5$, склад $D_5$ --- $B_5$ и
$B_7$. Наиболее загруженным является пункт $B_7$: в него поступает 22
единицы груза с трёх складов --- $D_1$, $D_2$ и $D_5$.

Полученное значение $Z_{\min}=611$ является минимально возможной суммарной
стоимостью перевозок. Любое изменение плана не может уменьшить суммарную
стоимость, поскольку все приведённые стоимости остаточных дуг
неотрицательны.

\clearpage

\section{Заключение}

В ходе выполнения лабораторной работы была решена транспортная задача с
промежуточными складами для варианта 4. Были выполнены следующие этапы.

\begin{enumerate}
    \item Изучена содержательная постановка задачи и проверены условия
          совместности: суммарный запас равен суммарной потребности (92),
          суммарная ёмкость складов достаточна (103).
    \item Построена математическая модель в виде задачи линейного
          программирования с ограничениями на запасы, заявки, сохранение
          груза на складах и ёмкости складов.
    \item Задача сведена к задаче о потоке минимальной стоимости в сети,
          содержащей 25 вершин и 83 дуги. Требуемая величина потока равна
          92.
    \item Методом последовательных кратчайших путей с потенциалами найден
          оптимальный поток. Вычисления потребовали 19 итераций; суммарная
          стоимость перевозок составила 611.
    \item Проведена проверка допустимости плана: все запасы вывезены, все
          заявки удовлетворены, груз на каждом складе сохранён, ёмкости
          складов не превышены.
    \item Оптимальность плана подтверждена с помощью потенциалов вершин:
          все приведённые стоимости остаточных дуг неотрицательны, поэтому
          в остаточной сети отсутствует отрицательный цикл.
\end{enumerate}

Полученный план перевозок обеспечивает минимальную суммарную стоимость
611 единиц. Склады $D_2$, $D_3$ и $D_4$ загружены полностью, а склады
$D_1$ и $D_5$ имеют свободный остаток. Результаты работы показывают, что
сведение двухэтапной транспортной задачи к задаче о потоке минимальной
стоимости позволяет эффективно учитывать ограничения на ёмкости складов и
находить точное оптимальное решение. Цель лабораторной работы достигнута.

\clearpage

\section{Список литературы}

\begin{enumerate}
    \item Канторович Л.В. Математические методы организации и планирования
          производства. --- Л.: Изд-во ЛГУ, 1939. --- 68 с.
    \item Форд Л.Р., Фалкерсон Д.Р. Потоки в сетях. --- М.: Мир, 1966. ---
          276 с.
    \item Беллман Р. Динамическое программирование. --- М.: Изд-во
          иностранной литературы, 1960. --- 400 с.
    \item Кормен Т., Лейзерсон Ч., Ривест Р., Штайн К. Алгоритмы:
          построение и анализ. --- М.: Вильямс, 2013. --- 1328 с.
    \item Пападимитриу Х., Стайглиц К. Комбинаторная оптимизация: алгоритмы
          и сложность. --- М.: Мир, 1985. --- 512 с.
    \item Лубенец Ю.В. Теория принятия решений: Учеб. пособие. --- Липецк:
          Изд-во ЛГТУ, 2011. --- 80 с.
\end{enumerate}
